\setcounter{section}{5}

\Needspace{5\baselineskip}
\hypertarget{ux89c6ux9891ux751fux6210ux6269ux6563ux6a21ux578bux7684ux4e0dux540cux76eeux6807ux51fdux6570}{%
\section{视频生成扩散模型的不同目标函数}\label{ux89c6ux9891ux751fux6210ux6269ux6563ux6a21ux578bux7684ux4e0dux540cux76eeux6807ux51fdux6570}}

\Needspace{5\baselineskip}
\hypertarget{ux4e00shortcut-forcingux6d41ux5339ux914dux6a21ux578bux7684ux6b65ux957fux62fcux63a5}{%
\subsection{一、Shortcut Forcing：流匹配模型的步长拼接}\label{ux4e00shortcut-forcingux6d41ux5339ux914dux6a21ux578bux7684ux6b65ux957fux62fcux63a5}}

\Needspace{5\baselineskip}
\hypertarget{ux4e00ux95eeux9898ux80ccux666f-3}{%
\subsubsection{（一）问题背景}\label{ux4e00ux95eeux9898ux80ccux666f-3}}

在视频生成或动力学建模中，传统扩散模型（Diffusion Models）或流匹配（Flow Matching）为了生成高质量的下一帧状态 \(x_{t+1}\)，通常需要数十步去噪迭代，例如 \(K=64\) 步。这在需要极快推理速度的强化学习交互中是不可接受的。

为了减少需要的步数，Dreamer 4 采用了 Shortcut Forcing，以提高大步长下的性能。

\Needspace{5\baselineskip}
\hypertarget{ux4e8cux6838ux5fc3ux601dux60f3-4}{%
\subsubsection{（二）核心思想}\label{ux4e8cux6838ux5fc3ux601dux60f3-4}}

模型训练时会随机抽取一个步长 \(d\)，例如 \(d\in\{1/64,2/64,4/64,\ldots,1\}\)。根据采样步长 \(d\) 选择不同损失函数，本质上是在构建一个从小步试探到大步跨越的自举蒸馏阶梯。

\Needspace{5\baselineskip}
当采样到最小步长 \(d=1/64\) 时，模型只需要进行一次很小的去噪，选择流匹配损失 \(\mathcal{L}_{\mathrm{flow}}\)。传统扩散模型大多采用 \(v\)-prediction，即预测速度方向：

\[
v=\epsilon-x
\]

其中 \(x\) 是干净数据，\(\epsilon\) 是高斯噪声。从频域（Fourier Transform）角度看，目标包含白噪声项。自然图像频谱通常近似遵循 \(1/f\) 规律，低频能量更高，而白噪声频谱是平坦的。强迫网络学习含有大量高频成分的 \(v\)，容易在长视频序列生成中导致误差累积（Error Accumulation）和高频伪影。

Shortcut Forcing 采用 \(x\)-prediction，即直接让网络预测干净的潜在表征。无论当前噪声时间步 \(\tau\in[0,1]\) 是多少，网络都尝试直接看透噪声，输出无噪的最终状态。为了让跨步预测保持一致，它引入 Bootstrap Loss：先将网络输出转换回 \(v\)-space，再通过 \((1-\tau)^2\) 的因子把 Bootstrap 损失缩放回 \(x\)-space 尺度进行优化。

\Needspace{5\baselineskip}
Dreamer 4 强制模型在任意噪声水平 \(\tau\in[0,1]\) 下，直接输出对最终干净表征的预测：

\[
\hat{x}_1=f_\theta(x_\tau,\tau,c)
\]

\Needspace{5\baselineskip}
基础流匹配损失在 \(x\) 空间内直接计算均方误差：

\[
\mathcal{L}_{\mathrm{flow}}=\left\|f_\theta(x_\tau,\tau,c)-x\right\|_2^2
\]

\Needspace{5\baselineskip}
当采样到较大步长 \(d=1/8\) 时，算法希望模型能够一步顶过好几步。如果直接让模型跨很大步长拟合真实数据，步长过大可能产生模糊或不符合物理规律的伪影。因此训练时使用自举一致性损失 \(\mathcal{L}_{\mathrm{bootstrap}}\)。工作流如下：

\begin{enumerate}
\def\labelenumi{\arabic{enumi}.}
\tightlist
\item
  算法先让模型用当前能力走两步较小的半步，即连续走两次 \(d/2=1/16\)。
\item
  将两次半步累加得到的结果作为临时教师答案（Teacher Target）。
\item
  再让模型直接跨越一个完整大步，即 \(d=1/8\)。
\item
  计算 \(\mathcal{L}_{\mathrm{bootstrap}}\)，强迫走一大步的输出必须等于走两小步的输出。
\end{enumerate}

\Needspace{5\baselineskip}
由于直接在 \(x\) 空间计算跨步损失尺度不一致，Dreamer 4 会把预测出的 \(\hat{x}_1\) 转换回 \(v\) 空间：

\[
\hat{v}_\tau=\frac{\hat{x}_1-x_\tau}{1-\tau}
\]

\Needspace{5\baselineskip}
计算两者的 MSE 后，再乘以 \((1-\tau)^2\) 的缩放因子，将损失重新拉平到 \(x\) 空间尺度：

\[
\mathcal{L}_{\mathrm{bootstrap}}=(1-\tau)^2\left\|\hat{v}_\tau-v_{\mathrm{target}}\right\|_2^2
\]

\Needspace{5\baselineskip}
为了让模型把有限参数容量集中在信息量最大的去噪阶段，也就是靠近干净数据的阶段，Dreamer 4 引入随信号水平线性递增的权重函数：

\[
w(\tau)=0.9\tau+0.1
\]

最终整体损失根据采样步长 \(d\) 选择 \(\mathcal{L}_{\mathrm{flow}}\) 或 \(\mathcal{L}_{\mathrm{bootstrap}}\)，并乘以权重 \(w(\tau)\) 进行优化。

\Needspace{5\baselineskip}
\hypertarget{ux4e8cedm-ux5f62ux5f0fux6269ux6563ux6a21ux578bux7684ux6e10ux53d8ux76eeux6807ux51fdux6570}{%
\subsection{二、EDM 形式：扩散模型的渐变目标函数}\label{ux4e8cedm-ux5f62ux5f0fux6269ux6563ux6a21ux578bux7684ux6e10ux53d8ux76eeux6807ux51fdux6570}}

\Needspace{5\baselineskip}
\hypertarget{ux4e00ux95eeux9898ux80ccux666f-4}{%
\subsubsection{（一）问题背景}\label{ux4e00ux95eeux9898ux80ccux666f-4}}

在标准 DDPM 中，神经网络被训练用来预测添加到干净图像上的噪声 \(\epsilon\)。在生成过程的起点，也就是时间步 \(T\)、噪声极大时，输入给网络的几乎是纯高斯噪声，此时信号被完全淹没。如果仍强迫网络预测输入中的噪声成分，网络会找到一个让损失很低的捷径：直接输出输入本身。数学上，这相当于网络退化成恒等映射。

扩散模型的本质是利用网络输出来估计数据分布的分数函数 \(\nabla_z\log p_T(z)\)。如果网络在采样第一步，也就是高噪声阶段，只输出输入本身，就会提供很差的分数估计。DDPM 可以依靠成百上千次微小去噪步逐渐修正，但在世界模型中，为保证智能体训练效率，通常只能负担少量去噪步，例如少于 10 步甚至 1 步。此时第一步误差无法被充分修正，会显著降低生成质量。

\Needspace{5\baselineskip}
\hypertarget{ux4e8cux6838ux5fc3ux7b97ux6cd5}{%
\subsubsection{（二）核心算法}\label{ux4e8cux6838ux5fc3ux7b97ux6cd5}}

\Needspace{5\baselineskip}
对于一个基于扩散模型的生成式世界模型，假设在 \(t+1\) 时刻初始噪声为 \(x_{t+1}\)，此前上下文为 \(\psi_t\)，可使用带跳连的自适应目标：

\[
D_\theta(x_{t+1}^{\tau},\psi_t)
=c_{\mathrm{skip}}^\tau x_{t+1}^{\tau}
+c_{\mathrm{out}}^\tau F_\theta(c_{\mathrm{in}}^\tau x_{t+1}^{\tau},\psi_t)
\]

\Needspace{5\baselineskip}
在这个公式下，神经网络 \(F_\theta\) 的实际训练目标变为：

\[
\mathcal{L}(\theta)=
\mathbb{E}\left[
\left\|
F_\theta(c_{\mathrm{in}}^\tau x_{t+1}^{\tau},\psi_t) -
\frac{1}{c_{\mathrm{out}}^\tau}
\left(x_{t+1}^{0}-c_{\mathrm{skip}}^\tau x_{t+1}^{\tau}\right)
\right\|_2^2
\right]
\]

\Needspace{5\baselineskip}
这个公式实现了一个渐变目标：

\begin{enumerate}
\def\labelenumi{\arabic{enumi}.}
\tightlist
\item
  当噪声极小时（\(\tau\to 0\)），系数 \(c_{\mathrm{skip}}\to 1\)。此时目标变成 \(x_{t+1}^0-x_{t+1}^{\tau}\)，即要求网络预测那一点点微小的残差噪声。
\item
  原理优势：如果此时仍然让网络预测干净图像 \(x^0\)，由于输入 \(x^\tau\) 已经和 \(x^0\) 几乎一样，网络输出与输入差异极小，会导致梯度消失，训练目标变得平庸。预测残差噪声则放大了这部分微小差异，迫使模型在最后阶段精细组织，完善高频细节。
\end{enumerate}

\Needspace{5\baselineskip}
\hypertarget{ux4e09ux6269ux6563-moe}{%
\subsubsection{（三）扩散 MoE}\label{ux4e09ux6269ux6563-moe}}

可采用双专家或多专家设计，以更好地实现自适应。高噪声专家负责早期去噪和全局结构，低噪声专家负责后期细节处理。

\FloatBarrier
\Needspace{5\baselineskip}
\hypertarget{ux53c2ux8003ux6587ux732e-143}{%
\subsection{参考文献}\label{ux53c2ux8003ux6587ux732e-143}}

\vspace{0.25em}
\begin{itemize}
\tightlist
\item
  Hafner, D., Yan, W., \& Lillicrap, T. (2025). \href{https://arxiv.org/abs/2509.24527}{Training Agents Inside of Scalable World Models}. arXiv:2509.24527.
\item
  Frans, K., Hafner, D., Levine, S., \& Abbeel, P. (2025). \href{https://arxiv.org/abs/2410.12557}{One Step Diffusion via Shortcut Models}. ICLR.
\item
  Karras, T., Aittala, M., Aila, T., \& Laine, S. (2022). \href{https://arxiv.org/abs/2206.00364}{Elucidating the Design Space of Diffusion-Based Generative Models}. NeurIPS.
\end{itemize}

\clearpage
